% Gate and labels. DRAFT, 2026-11-24. Needs booktabs and amsmath.
% Every number here is a DEVELOPMENT result: surrogate VLM, validation seeds, scenario not frozen.
% Replace Section "Results" with the final rerun (frozen scenario, real VLM, test seeds) before submission.
% Sources: AGENT_LOG.md 2026-11-16 to 2026-11-24, results/gate_val.json, results/gate_holdout_*, results/benchmark.md.
\providecommand{\todo}[1]{\textbf{[TODO: #1]}}

\section{The gate: when is it worth asking?}
\label{sec:gate}

The gate reads the fast stack's output on every frame and decides whether to pause the drone and ask the VLM one
yes/no question. It has two parts that we can choose separately: what it looks at (its features), and what it is
trained to predict (its labels). We keep the features and the model fixed and change only the labels.

\subsection{Labels}
\label{sec:gate-labels}

Let $d^{\mathrm{gt}}$ and $\hat d$ be the nearest distance in the forward corridor (the region of the image the
drone is about to fly through), from simulator ground truth and from the estimated depth, and let
$d_{\mathrm{safe}} = 8$\,m be the distance at which the controller starts to avoid. Ground truth is used only to
build labels, never by the controller, the gate or any policy in flight.

\paragraph{L1, missed obstacle.} A frame is positive when the estimate says the way is clear and it is not:
\[
  y_{\mathrm{L1}} = 1 \iff d^{\mathrm{gt}} < d_{\mathrm{safe}} \le \hat d .
\]
L1 says that perception is wrong, not that asking would help.

\paragraph{L3, worth asking.} L3 asks about the outcome. We fly each seed once without asking. At every second
frame whose flight crashes within the horizon $H = 3$\,s, we restore the exact flight state and fly on from there,
asking exactly once at that frame: the drone brakes and holds for $d = 1$\,s of simulated time, the VLM answers,
and a ``yes'' makes the drone avoid. The frame is positive only if the flight that asked does not crash:
\[
  y_{\mathrm{L3}} = 1 \iff \text{crash without asking} \ \wedge\ \text{no crash when asking.}
\]
The asking branch gets $H + d$ seconds, so both branches fly for the same time after the question; otherwise the
pause alone would push a crash past the horizon and look like a save. The simulator steps poses directly, so
restoring a state does not need a simulator snapshot. We do not label as positive a frame where asking
\emph{causes} a crash, because that would train the gate to ask where asking hurts. Only frames whose no-ask flight
crashes within $H$ are branched, so this case does not occur in our labels and is not measured.

L1 and L3 differ in an informative way: on the 20 train flights we compared first, they agree on 346 of 495 frames.
\todo{Final agreement figure from the rerun.}

\subsection{Features and model}
\label{sec:gate-model}

Both gates use the same 17 features, all available in flight: seven statistics of the estimated depth in the corridor
(minimum, 5th and 25th percentiles, median, variance, mean gradient, share of far pixels); four changes of those
statistics per metre flown, so they mean the same at any speed; and six image features of the RGB corridor crop that
target thin structures (gradient, edge density, horizontal and vertical line response, contrast, brightness). The model
is a small multilayer perceptron (hidden layers of 32 and 16 units, standardized inputs, $\ell_2$ penalty $10^{-3}$,
early stopping, fixed seed). Training and validation flights are separated by seed, never by frame.

In flight, the gate asks when its score exceeds a cut $\tau$, at most once per second. Cuts are set on validation
frames so that the gate asks on 5\% or 10\% of them; they are never set on test seeds.

\subsection{Baselines}
\label{sec:gate-baselines}

We compare against a policy that never asks, one that asks at random with a probability chosen so that its calls per
flight are close to the gate's, a threshold on the estimated corridor minimum (asking when it is \emph{high}, the
direction that predicted L3 best on train; its cut is set on validation like the gates'), an oracle that asks exactly
when L1 is positive, and a policy that asks as often as it can.

\section{Results}
\label{sec:gate-results}

\noindent\textbf{Read this first.} Unless stated otherwise, the numbers below use a \emph{surrogate VLM}: a random
answerer with the recall (0.80) and false-alarm rate (0.14) we measured for Qwen3-VL-8B on frames where depth misses
the obstacle. They come from validation seeds and a scenario that is not frozen. They are development results and
will be replaced. \todo{Rerun on the frozen scenario with real VLM answers, then test seeds once.}

\subsection{Gate-L1 on the frame bank}

On 1{,}830 validation frames from a held-out set of bank episodes (281 positive, 19 flights), the L1 gate reaches
AUROC 0.829 [0.753, 0.906] and AUPRC 0.566 (base rate 0.154). Depth variance alone reaches 0.425, the best single
depth feature 0.695, and the gate without image features 0.635. Most of the gain comes from the image features.
Holding out one obstacle family from training keeps most of the quality for thin obstacles (AUROC against all
negatives for the held-out family: poles 0.81, 3\,cm wires 0.71) and loses it for glass (0.04, with only two
positive validation flights). A value below 0.5 means the gate ranked glass as \emph{safer} than clear frames. We
report this as a limitation: the gate learned glass only from glass.

\subsection{Gate-L3 against gate-L1 on replay frames}

Both gates were trained on the same frames (100 flights, 2{,}961 frames, 62 of 100 crashing, 724 L3-positive) and
scored on 60 flights (1{,}999 frames, 28 crashing, 346 L3-positive). Table~\ref{tab:offline} compares them at equal
query rates, with intervals from resampling whole flights.

\begin{table}[t]
\centering
\caption{Gate-L3 minus gate-L1 on validation replay frames (surrogate VLM). L3 recall: share of worth-asking
frames on which the gate asks. Wasted: share of asks that change nothing.}
\label{tab:offline}
\begin{tabular}{@{}lccc@{}}
\toprule
 & 5\% & 10\% & 20\% \\
\midrule
L3 recall, difference & $+0.058$ & $+0.101$ & $+0.197$ \\
\quad 95\% interval & $[0.015, 0.122]$ & $[0.025, 0.198]$ & $[0.089, 0.294]$ \\
\midrule
AUROC against L3 & \multicolumn{3}{c}{gate-L3 0.857 [0.809, 0.900]; gate-L1 0.717 [0.651, 0.778]} \\
Wasted asks at 5\% & \multicolumn{3}{c}{gate-L3 0.290; gate-L1 0.490; random 0.880} \\
\bottomrule
\end{tabular}
\end{table}

Training on L3 makes the gate rank worth-asking frames better and waste fewer asks. It does not show a gain in the
share of crashing flights saved: every interval for that difference includes zero. With about 12 savable frames per
savable flight, even random asking at 10\% catches one on 61\% of them, so this measure saturates.

\subsection{In flight}

Table~\ref{tab:flight} flies the same 60 validation seeds with each policy, where each ask costs one second of
braking and changes what follows.

\begin{table}[t]
\centering
\caption{Closed loop, surrogate VLM, validation seeds, 60 flights per policy. Wilson 95\% intervals for the crash rate.
Development result; the scenario is not frozen.}
\label{tab:flight}
\begin{tabular}{@{}lrlrr@{}}
\toprule
Policy & Crashes & Crash rate & Calls/flight & Mission (s) \\
\midrule
never                  & 29 & 0.48 [0.36, 0.61] & 0.0  & 6.6 \\
\midrule
gate-L3, 5\%           & 13 & 0.22 [0.13, 0.34] & 0.6  & 9.1 \\
gate-L1, 5\%           & 11 & 0.18 [0.11, 0.30] & 0.8  & 9.6 \\
random, $p=0.01$       & 22 & 0.37 [0.26, 0.49] & 0.8  & 8.6 \\
threshold, 5\%         & 23 & 0.38 [0.27, 0.51] & 0.5  & 8.1 \\
\midrule
gate-L3, 10\%          &  7 & 0.12 [0.06, 0.22] & 1.2  & 10.6 \\
gate-L1, 10\%          &  5 & 0.08 [0.04, 0.18] & 1.2  & 10.9 \\
random, $p=0.02$       & 17 & 0.28 [0.19, 0.41] & 1.3  & 9.7 \\
threshold, 10\%        & 17 & 0.28 [0.19, 0.41] & 1.0  & 10.0 \\
\midrule
random, $p=0.10$       &  4 & 0.07 [0.03, 0.16] & 5.6  & 16.3 \\
oracle                 &  0 & 0.00 [0.00, 0.06] & 5.5  & 16.0 \\
always                 &  2 & 0.03 [0.01, 0.11] & 54.2 & 59.6 \\
\bottomrule
\end{tabular}
\end{table}

\emph{Gates against baselines.} At similar calls per flight, each gate crashes less than random and less than the
threshold in all eight paired comparisons (crash-rate differences of 15 to 20 points; every 95\% bootstrap interval
excludes zero). After a Holm correction for the eight McNemar tests, four remain at the 0.05 level. The intervals
for the difference in calls per flight include zero, so calls are not shown to differ, which is weaker than showing
that they are equal.

\emph{L3 against L1.} No closed-loop difference is visible: gate-L3 crashes on two more flights than gate-L1 at each
cut (McNemar $p = 0.80$ and $0.75$). Detecting a gap this small would need roughly 400 or more flights. The offline
advantage of L3 therefore does not carry over in this experiment.

\emph{Why the remaining crashes happen.} In the baselines, nearly all crashes are missed asks (21 of 22 for random at
$p=0.01$, 23 of 23 for the 5\% threshold). At the 10\% cut, four crashes per gate follow a ``no'' from the VLM,
which no gate can remove. Asking always is not a usable reference here: it spends the whole 60\,s limit asking and
completes no mission.

\subsection{Does it hold with the real VLM?}

A first check with Qwen3-VL-8B in the loop (30 flights, five seeds, 25 asks outside the oracle) shows that the model
runs in the flight job at about 0.15\,s per call, but that its decision cut does not transfer. With the true
depth beyond 30\,m and nothing placed within 30\,m, it answered ``yes'' on 12 of 13 asks at the cut tuned on the frame
bank, where the same cut gives 10\% false alarms. Under 15\,m it answered ``yes'' on 6 of 7. Confidence still
separates near from far, so this is a calibration problem. Two likely causes are that the bank images we inspected
show one location while flights use five, and that the prompt asks whether an object \emph{crosses} the region
without saying how near. Until the cut is retuned on train flights, the surrogate's 0.14 false-alarm rate is
too optimistic for these flights. \todo{Retune (SB), rerun labels, gates and sweep.}

\subsection{Limitations}
\label{sec:gate-limits}
\begin{itemize}
\item \textbf{Surrogate answers.} Errors are independent draws; a real VLM tends to miss the same object every time.
\item \textbf{Tuned on these seeds.} Gate and threshold cuts were set on the same 60 validation seeds,
      which favours them slightly over random. The first 432 flights also used an earlier judge seeding scheme.
\item \textbf{Start poses.} The map itself obstructs climbing drones at three of five start poses. Of 188 crashes over all
      policies, 13 hit scenery that we did not place, all at two poses; none came from never, random at 1--2\%, or the tuned threshold.
\item \textbf{Difficulty.} The never-ask policy crashes on 48\% of these flights, above our 15 to 30\% target.
\item \textbf{Simulator noise.} Two of 60 never-ask seeds ended differently on a second run (one crash outcome, one
      crash time shifted by 0.1\,s); two of 233 flights that made no calls ended unlike the never-ask flight.
\item \textbf{Small counts.} Per-family results have single-digit counts and are not used for claims.
\end{itemize}
